% Source: arXiv:2203.12563v3, REsubmission.tex, lines 1695--1777.

\subsection{Joint compression at the two boundary sites}
\label{sec:mpo_joint_edge_normalization}

Retain $C_x$, $V_x$, $L=U_LP_L$, and $R=U_RP_R$ from
Section~\ref{sec:mpo_joint_boundary_columns}. Write $D_x=D_x^0+D_x^1$
and let $\iota_0:\Fin{d_0}\to\mathcal I$ denote the shared first physical
alphabet. The spaces $\C^{\mathcal B_L}$ and $\C^{\mathcal B_R}$ include
all block labels jointly.

\begin{definition}[One-sided supports with exterior multiplicities]
    \label{def:mpo_joint_edge_cores}
    \lean{MPSTensor.MPOSymmetry.jointEndpointFirstEdgeCoreMap,
        MPSTensor.MPOSymmetry.jointEndpointLastEdgeCoreMap}
    \leanok
    For arbitrary nonnegative exterior dimensions $E_x$, set
    $\mathcal B_L(E)=\coprod_x(\Fin{E_x}\times\Fin{D_x^0})$ and
    $\mathcal B_R(E)=\coprod_x(\Fin{D_x^0}\times\Fin{E_x})$.
    Define two linear maps by
    \begin{align}
        \Phi_{L,E}(Y)_{(x,a,b),j} & =((A_x^0)^jY_x)_{ba},
                                  & Y_x                   & \in\C^{D_x^0\times E_x}, \\
        \Phi_{R,E}(Z)_{i,(x,c,e)} & =(Z_x(A_x^0)^i)_{ec},
                                  & Z_x                   & \in\C^{E_x\times D_x^0}.
    \end{align}
    Their ranges are the respective one-sided core supports.
\end{definition}

\begin{definition}[Actual cropped edge coefficients]
    \label{def:mpo_joint_cropped_edges}
    \lean{MPSTensor.MPOSymmetry.jointMixedFirstEdgeBoundaryMap,
        MPSTensor.MPOSymmetry.jointMixedLastEdgeBoundaryMap}
    \leanok
    \uses{def:mpo_joint_mixed_family}
    For $X_x\in\MN{D_x}$, define
    \begin{align}
        B_L(X)_{i,j} & =\sum_x\tr(C_x^iV_xV_x^*C_x^{\iota_0(j)}X_x), \\
        B_R(X)_{i,j} & =\sum_x\tr(C_x^{\iota_0(i)}V_xV_x^*C_x^jX_x).
    \end{align}
    Only the adjacent bulk site is restricted to the shared alphabet;
    it retains all $d_0$ coordinates.
\end{definition}

\begin{theorem}[Shared letters at the virtual inclusion]
    \label{thm:mpo_joint_shared_letter_inclusion}
    \lean{MPSTensor.MPOSymmetry.jointMixedEndpointBase_firstPhysical_mul_inclusion,
        MPSTensor.MPOSymmetry.jointMixedEndpointBase_inclusion_adjoint_mul_firstPhysical}
    \leanok
    \uses{def:mpo_joint_mixed_family}
    Every shared first-endpoint letter satisfies
    $C_x^{\iota_0(i)}V_x=V_x(A_x^0)^i$ and
    $V_x^*C_x^{\iota_0(i)}=(A_x^0)^iV_x^*$.
\end{theorem}
\begin{proof}
    \leanok
    \uses{def:mpo_joint_mixed_family}
    In the virtual sum $\C^{D_x^0}\oplus\C^{D_x^1}$, the letter
    $C_x^{\iota_0(i)}$ is $\operatorname{diag}((A_x^0)^i,0)$.
    Multiplication by the first coordinate inclusion gives both identities.
\end{proof}

\begin{theorem}[Actual coefficients factor through one-sided cores]
    \label{thm:mpo_joint_edge_factors}
    \lean{MPSTensor.MPOSymmetry.jointMixedFirstEdgeBoundaryMap_eq_columns_core,
        MPSTensor.MPOSymmetry.jointMixedLastEdgeBoundaryMap_eq_columns_core}
    \leanok
    \uses{def:mpo_joint_cropped_edges,def:mpo_joint_edge_cores,def:mpo_joint_boundary_columns}
    With $E_x=D_x$, the actual maps satisfy
    \begin{align}
        B_L(X) & =(L\otimes I_{d_0})\Phi_{L,D}((V_x^*X_x)_x),
        \label{eq:mpo_joint_first_core_factor}
    \end{align}
    \begin{align}
        B_R(X) & =(I_{d_0}\otimes R)\Phi_{R,D}((X_xV_x)_x).
        \label{eq:mpo_joint_last_core_factor}
    \end{align}
    No spanning or positive-dimension hypothesis is required.
\end{theorem}
\begin{proof}
    \leanok
    \uses{thm:mpo_joint_shared_letter_inclusion,def:mpo_joint_boundary_columns}
    Substitute the virtual inclusion identities and expand the trace:
    \begin{align}
        B_L(X)_{i,j}
         & =\sum_{x,a,b}L_{i;(x,a,b)}
        ((A_x^0)^jV_x^*X_x)_{ba},     \\
        B_R(X)_{i,j}
         & =\sum_{x,c,e}R_{j;(x,c,e)}
        (X_xV_x(A_x^0)^i)_{ec}.
    \end{align}
    The second equality uses cyclicity of the trace. These are
    (\ref{eq:mpo_joint_first_core_factor}) and
    (\ref{eq:mpo_joint_last_core_factor}) in coordinates.
\end{proof}

\begin{theorem}[Every rectangular boundary is attained]
    \label{thm:mpo_joint_boundary_extraction}
    \lean{MPSTensor.MPOSymmetry.jointMixedFirstBoundaryExtraction_surjective,
        MPSTensor.MPOSymmetry.jointMixedLastBoundaryExtraction_surjective}
    \leanok
    The maps $(X_x)_x\mapsto(V_x^*X_x)_x$ and
    $(X_x)_x\mapsto(X_xV_x)_x$ are surjective onto their respective
    rectangular boundary spaces. The assertion includes empty label
    sets and zero-dimensional fibers.
\end{theorem}
\begin{proof}
    \leanok
    Given $Y_x\in\C^{D_x^0\times D_x}$, choose $X_x=V_xY_x$.
    Given $Z_x\in\C^{D_x\times D_x^0}$, choose $X_x=Z_xV_x^*$.
    Both claims follow from $V_x^*V_x=I$.
\end{proof}

\begin{definition}[Compression into joint polar coordinates]
    \label{def:mpo_joint_polar_edge_compression}
    \lean{MPSTensor.MPOSymmetry.jointMixedFirstEdgeCompressedMap,
        MPSTensor.MPOSymmetry.jointMixedLastEdgeCompressedMap}
    \leanok
    \uses{def:mpo_joint_cropped_edges,thm:mpo_joint_boundary_polar}
    Define $\widehat B_L=(U_L^*\otimes I_{d_0})B_L$ and
    $\widehat B_R=(I_{d_0}\otimes U_R^*)B_R$.
\end{definition}

\begin{theorem}[Only the boundary polar weights remain]
    \label{thm:mpo_joint_polar_edge_coefficients}
    \lean{MPSTensor.MPOSymmetry.jointMixedFirstEdgeCompressedMap_eq_polarPos_core,
        MPSTensor.MPOSymmetry.jointMixedLastEdgeCompressedMap_eq_polarPos_core,
        MPSTensor.MPOSymmetry.jointMixedFirstEdge_normalize_compressed,
        MPSTensor.MPOSymmetry.jointMixedLastEdge_normalize_compressed}
    \leanok
    \uses{def:mpo_joint_polar_edge_compression,def:mpo_joint_edge_cores,
        thm:mpo_joint_boundary_polar}
    Suppose both endpoint families span their block algebras simultaneously
    at one site. Then
    \begin{align}
        \widehat B_L(X) & =(P_L\otimes I_{d_0})\Phi_{L,D}((V_x^*X_x)_x), \\
        \widehat B_R(X) & =(I_{d_0}\otimes P_R)\Phi_{R,D}((X_xV_x)_x),   \\
        (P_L^{-1}\otimes I_{d_0})\widehat B_L(X)
                        & =\Phi_{L,D}((V_x^*X_x)_x),                     \\
        (I_{d_0}\otimes P_R^{-1})\widehat B_R(X)
                        & =\Phi_{R,D}((X_xV_x)_x).
    \end{align}
    The matrices $P_L,P_R$ retain all cross-label entries.
\end{theorem}
\begin{proof}
    \leanok
    \uses{thm:mpo_joint_boundary_polar,thm:mpo_joint_edge_factors}
    Since $L=U_LP_L$ and $U_L^*U_L=I$, one has $U_L^*L=P_L$;
    likewise $U_R^*R=P_R$. Apply these identities to
    (\ref{eq:mpo_joint_first_core_factor}) and
    (\ref{eq:mpo_joint_last_core_factor}). Positive definiteness makes
    the two weights invertible, so multiplying by their inverses gives
    the normalized coefficients.
\end{proof}

\begin{theorem}[Exact normalized edge supports]
    \label{thm:mpo_joint_normalized_edge_supports}
    \lean{MPSTensor.MPOSymmetry.map_jointMixedFirstEdgeCompressed_range_eq_core,
        MPSTensor.MPOSymmetry.map_jointMixedLastEdgeCompressed_range_eq_core}
    \leanok
    \uses{def:mpo_joint_polar_edge_compression,def:mpo_joint_edge_cores}
    Under the simultaneous spanning hypotheses,
    \begin{align}
        (P_L^{-1}\otimes I_{d_0})\operatorname{ran}\widehat B_L
         & =\operatorname{ran}\Phi_{L,D}, \\
        (I_{d_0}\otimes P_R^{-1})\operatorname{ran}\widehat B_R
         & =\operatorname{ran}\Phi_{R,D}.
    \end{align}
    Neither equality changes an interior tensor or assumes that the
    polar frame fills its ambient physical alphabet.
\end{theorem}
\begin{proof}
    \leanok
    \uses{thm:mpo_joint_polar_edge_coefficients,thm:mpo_joint_boundary_extraction}
    The coefficient identities give the forward inclusions.
    The reverse inclusions follow by choosing the padded square
    boundaries from Theorem~\ref{thm:mpo_joint_boundary_extraction}.
\end{proof}

For a chain of length at least three, the first and last constraints
act on distinct edges. At length two they are descriptions of a single
edge, whose kernel must be transformed by both boundary changes at once.
No sum of edge operators is asserted here.

\begin{definition}[Euclidean boundary normalization and support]
    \label{def:mpo_joint_edge_normalization_equiv}
    \lean{MPSTensor.MPOSymmetry.jointMixedFirstEdgeNormalizationEquivES,
        MPSTensor.MPOSymmetry.jointMixedLastEdgeNormalizationEquivES,
        MPSTensor.MPOSymmetry.jointMixedFirstEdgeCompressedSupportES,
        MPSTensor.MPOSymmetry.jointMixedLastEdgeCompressedSupportES,
        MPSTensor.MPOSymmetry.jointEndpointFirstEdgeCoreSupportES,
        MPSTensor.MPOSymmetry.jointEndpointLastEdgeCoreSupportES}
    \leanok
    \uses{def:mpo_joint_edge_cores,def:mpo_joint_polar_edge_compression,
        thm:mpo_joint_boundary_polar}
    Equip all coefficient spaces with their Euclidean inner products.
    Set $S_L=\operatorname{ran}\widehat B_L$,
    $S_R=\operatorname{ran}\widehat B_R$,
    $T_{L,E}=\operatorname{ran}\Phi_{L,E}$, and
    $T_{R,E}=\operatorname{ran}\Phi_{R,E}$.
    Under simultaneous one-site spanning, define the invertible maps
    $e_L=P_L^{-1}\otimes I_{d_0}$ and
    $e_R=I_{d_0}\otimes P_R^{-1}$. Their inverses are
    $P_L\otimes I_{d_0}$ and $I_{d_0}\otimes P_R$.
\end{definition}

\begin{theorem}[Normalized Euclidean support]
    \label{thm:mpo_joint_edge_euclidean_support}
    \lean{MPSTensor.MPOSymmetry.map_jointMixedFirstEdgeNormalization_support,
        MPSTensor.MPOSymmetry.map_jointMixedLastEdgeNormalization_support}
    \leanok
    \uses{def:mpo_joint_edge_normalization_equiv}
    With $E_x=D_x$, one has $e_L(S_L)=T_{L,D}$ and $e_R(S_R)=T_{R,D}$.
\end{theorem}
\begin{proof}
    \leanok
    \uses{thm:mpo_joint_normalized_edge_supports}
    The underlying coefficient vectors are unchanged when their
    Euclidean norms are specified, so the two range equalities apply directly.
\end{proof}

\begin{theorem}[Canonical constraints in both directions]
    \label{thm:mpo_joint_edge_deformed_constraints}
    \lean{MPSTensor.MPOSymmetry.jointMixedFirstEdgeNormalization_deformed_constraints,
        MPSTensor.MPOSymmetry.jointMixedLastEdgeNormalization_deformed_constraints}
    \leanok
    \uses{def:mpo_joint_edge_normalization_equiv}
    For either boundary, let $e(S)=T$ be the corresponding normalization
    and write $h_S=\Pi_{S^\perp}$, $h_T=\Pi_{T^\perp}$.
    For an invertible map $f$, define its canonical deformation by
    $\mathcal D_f(h)=\Pi_{(f^{-1}(\ker h))^\perp}$. Then
    $\mathcal D_{e^{-1}}(h_S)=h_T$ and
    $\mathcal D_e(h_T)=h_S$.
\end{theorem}
\begin{proof}
    \leanok
    \uses{thm:mpo_joint_edge_euclidean_support}
    Since $\ker h_S=S$ and $\ker h_T=T$, the two pulled-back kernels
    are $e(S)=T$ and $e^{-1}(T)=S$. Orthogonal projection onto their
    complements gives the asserted equalities.
\end{proof}

\begin{definition}[Isometric compression]
    \label{def:mpo_rectangular_compression}
    \lean{LinearIsometry.compression}
    \leanok
    For finite-dimensional complex Hilbert spaces $E,F$, an isometry
    $U:E\to F$, and an operator $H:F\to F$, define $K=U^*HU$ on $E$.
\end{definition}

\begin{theorem}[Compression intertwines on a reducing image]
    \label{thm:mpo_rectangular_compression_intertwines}
    \lean{LinearIsometry.apply_compression_of_commute,
        LinearIsometry.compression_adjoint_apply_of_commute,
        LinearIsometry.ker_compression_of_commute,
        LinearIsometry.ker_compression_eq_map_adjoint_of_commute,
        LinearIsometry.map_mem_orthogonal_ker_compression_of_commute}
    \leanok
    \uses{def:mpo_rectangular_compression}
    Suppose that $UU^*H=HUU^*$. Then
    \begin{align}
        UK                & =HU,                         & KU^* & =U^*H, \\
        \ker K            & =U^{-1}(\ker H)=U^*(\ker H), &
        U((\ker K)^\perp) & \subseteq(\ker H)^\perp.
    \end{align}
    No surjectivity or self-adjointness assumption is needed.
\end{theorem}
\begin{proof}
    \leanok
    \uses{def:mpo_rectangular_compression}
    Use $U^*U=I$ and the commutation hypothesis to obtain
    $UK=UU^*HU=HU$ and $KU^*=U^*HUU^*=U^*H$.
    Injectivity of $U$ gives $\ker K=U^{-1}(\ker H)$.
    The second intertwining identity sends $\ker H$ into $\ker K$,
    while $v=U^*(Uv)$ gives the reverse inclusion in $U^*(\ker H)$.
    If $v\perp\ker K$ and $w\in\ker H$, then
    $\langle w,Uv\rangle=\langle U^*w,v\rangle=0$.
\end{proof}

\begin{theorem}[A norm gap restricts without loss]
    \label{thm:mpo_rectangular_compression_gap}
    \lean{LinearIsometry.norm_gap_compression_of_commute}
    \leanok
    \uses{def:mpo_rectangular_compression}
    Suppose that $UU^*H=HUU^*$ and that $\delta\in\R$ satisfies
    $\delta\lVert w\rVert\le\lVert Hw\rVert$ for every
    $w\in(\ker H)^\perp$. Then
    $\delta\lVert v\rVert\le\lVert Kv\rVert$ for every
    $v\in(\ker K)^\perp$. The assertion includes zero-dimensional
    domain or codomain.
\end{theorem}
\begin{proof}
    \leanok
    \uses{thm:mpo_rectangular_compression_intertwines}
    For $v\perp\ker K$, apply the given estimate to $Uv$ and use
    $HUv=UKv$. Since $U$ preserves norms,
    $\delta\lVert v\rVert=\delta\lVert Uv\rVert
        \le\lVert HUv\rVert=\lVert Kv\rVert$.
\end{proof}
